\SourceLeaf{18}
whence $r=\BareRoot{((a-x)^2+(b-y)^2+(c-z)^2)}$
\[
\begin{aligned}
\cos MX&=\frac{a-x}{r}\\
\cos MY&=\frac{b-y}{r}\\
\cos MZ&=\frac{c-z}{r}
\end{aligned}
\]
Finally, by a well-known theorem, we have
\[
\cos MQ=\cos MX.\cos QX+\cos MY.\cos QY+\cos MZ.\cos QZ
\]
or
\[
\cos MQ=\pm\frac{T(a-x)+U(b-y)+V(c-z)}{r\BareRoot{(TT+UU+VV)}}
\]

\Article{10}
\par
To carry out the integration of the differential expressions over the whole surface, we must now transform them so that they contain only two variables. This can be done by eliminating one of the variables $x,y,z$ using the equation $W=0$, though this usually leads to formulas that are less easy to handle. It is therefore preferable to introduce two new indeterminates $p,q$, so that each of $x$, $y$, and $z$ is regarded as a function of these indeterminates.

\begingroup\rightskip=0pt plus 1em
As soon as definite values are assigned to $p,q$, $x,y,z$ will also be determined. That is, a definite point on the surface of the body will correspond to those values. This correspondence in both directions becomes clearer if we imagine an unbounded plane whose points are given by rectangular coordinates $p,q$. Every point of the plane will then correspond to a point on the surface of the body, and to just one point if the correspondence is arranged so that $x,y,z$ are single-valued functions of the indeterminates $p,q$. Conversely, if $x,y,z$ also determine $p$ and $q$ fully and unambiguously, then clearly there will be just one point of the plane corresponding to each point on the surface of the body. In this case the plane must extend to infinity in every direction to cover the entire surface of the body. Otherwise we need consider only a portion of the plane, with finite or infinite bounds, that represents, as it were, the surface of the body. Imagine that the plane is divided by infinitely many straight lines, some parallel to the axis of abscissas and others perpendicular to it,
